\begin{table}[H]
\centering
\caption{Kebutuhan Nonfungsional (NFR) Sistem}\label{tbl:kebutuhan-nonfungsional}
\renewcommand{\arraystretch}{1.3}
\begin{tabular}{|l|p{0.8\textwidth}|}
\hline
\textbf{Kode} & \textbf{Deskripsi Kebutuhan Nonfungsional} \\
\hline
NFR-01 & \textbf{(Usability)} Antarmuka interaksi sosial harus terintegrasi secara intuitif dalam platform yang sama dengan platform pembelajaran, meminimalkan beban kognitif siswa saat berpindah antara belajar mandiri dan berkolaborasi. \\
\hline
NFR-02 & \textbf{(Interpretability)} Sistem harus mampu memberikan indikasi alasan rekomendasi (baik materi maupun interaksi sosial). \\
\hline
NFR-03 & \textbf{(Performance)} Pembaruan model pengguna (baik skor individu maupun sosial) harus terjadi dalam waktu nyata (\textit{near real-time}) setelah interaksi selesai dilakukan. \\
\hline
NFR-04 & \textbf{(Scalability)} Sistem harus mampu menangani konkurensi interaksi sosial dari satu kelas (±50 pengguna) secara simultan. \\
\hline
NFR-05 & \textbf{(Privacy)} Mekanisme interaksi sosial harus dirancang sedemikian rupa sehingga data privasi spesifik (seperti nilai asesmen mentah) tidak terekspos ke pengguna lain. \\
\hline
NFR-06 & \textbf{(Accessibility)} Sistem dapat diakses melalui peramban web standar. \\
\hline
\end{tabular}
\end{table}